% GENERATED FILE. Edit canonical lessons, then run npm run generate:books.
% canonical-source-hash: fnv1a64:d4aae1c592ed2be2
% canonical-lessons: LA-C399-perfidia, LA-C399-perfidus, LA-C399-fallacia, LA-C399-decipio, LA-C399-testor

\chapter{Trust and Betrayal}
\label{ch:la-trust-and-betrayal}
\hlchaptermodality{399}

\begin{chapteropening}
\textbf{By the end of this chapter:} \emph{I can talk about bad faith and a treacherous friend, a trick and a deception, and call the gods to witness.}

Complete the last lesson of chapter 399: I can talk about bad faith and a treacherous friend, a trick and a deception, and call the gods to witness.
\end{chapteropening}

\section[perfidia, perfidiae]{perfidia, perfidiae — bad faith, treachery}

\label{lesson:LA-C399-perfidia}



\begin{quote}
Before the new one: say the Latin for fixed, then the Latin for void.
\end{quote}

\subsection*{You'll want to know: perfidia, perfidiae}
\textbf{perfidia, perfidiae} — \textquotedblleft{}bad faith, treachery\textquotedblright{}.

\begin{quote}
Perfidiam eius omnēs nōvērunt.
\end{quote}

\begin{quote}
\emph{Everyone knew of his treachery.}
\end{quote}

\emph{perfidia, perfidiae} f., \emph{per-} past + \emph{fidēs}, faith: going past one's word. English \emph{perfidy}.

\begin{cousinweb}
\emph{perfidia} is built on \emph{perfidus} \textquotedblleft{}faithless\textquotedblright{} (Lewis \& Short), from \emph{per-} \textquotedblleft{}past, beyond\textquotedblright{} and \emph{fidēs} \textquotedblleft{}faith\textquotedblright{}: going beyond one's word. English \emph{perfidy} comes from it.
\end{cousinweb}

\begin{grammarlens}[title={What you've built}]
One more word for telling, explaining and describing at B1.
\end{grammarlens}

\subsection*{Your turn}
\begin{itemize}
\raggedright
  \item \emph{Say it:} \emph{perfidia, perfidiae}
  \item \emph{Say it:} \emph{perfidia, perfidiae}, once more
  \item \emph{Say it:} a sentence of your own with \emph{perfidia, perfidiae}
  \item \emph{Recall:} say the Latin for fixed, then the Latin for void, then say \emph{perfidia, perfidiae} again
\end{itemize}

\begin{tcolorbox}[breakable,colback=teal!4,colframe=teal!35!black,title={Before you move on}]
What does perfidia, perfidiae mean? (\textquotedblleft{}Bad faith, treachery\textquotedblright{}.) Say it once more.
\end{tcolorbox}
\section[perfidus, perfida, perfidum]{perfidus, perfida, perfidum — treacherous, faithless}

\label{lesson:LA-C399-perfidus}



\begin{quote}
Before the new one: say the Latin for void, then the Latin for bad faith.
\end{quote}

\subsection*{You'll want to know: perfidus, perfida, perfidum}
\textbf{perfidus, perfida, perfidum} — \textquotedblleft{}treacherous, faithless\textquotedblright{}.

\begin{quote}
Perfidō amīcō nēmō crēdit.
\end{quote}

\begin{quote}
\emph{Nobody trusts a faithless friend.}
\end{quote}

\emph{perfidus, perfida, perfidum}: the opposite of \emph{fīdus}, faithful.

\begin{cousinweb}
\emph{perfidus} is built from \emph{per-} \textquotedblleft{}past, beyond\textquotedblright{} and \emph{fidēs} \textquotedblleft{}faith\textquotedblright{} (EtymDB; Lewis \& Short): one who goes beyond his word.
\end{cousinweb}

\begin{grammarlens}[title={What you've built}]
One more word for telling, explaining and describing at B1.
\end{grammarlens}

\subsection*{Your turn}
\begin{itemize}
\raggedright
  \item \emph{Say it:} \emph{perfidus, perfida, perfidum}
  \item \emph{Say it:} \emph{perfidus, perfida, perfidum}, once more
  \item \emph{Say it:} a sentence of your own with \emph{perfidus, perfida, perfidum}
  \item \emph{Recall:} say the Latin for void, then the Latin for bad faith, then say \emph{perfidus, perfida, perfidum} again
\end{itemize}

\begin{tcolorbox}[breakable,colback=teal!4,colframe=teal!35!black,title={Before you move on}]
What does perfidus, perfida, perfidum mean? (\textquotedblleft{}Treacherous, faithless\textquotedblright{}.) Say it once more.
\end{tcolorbox}
\section[fallācia, fallāciae]{fallācia, fallāciae — a deceit, a trick}

\label{lesson:LA-C399-fallacia}



\begin{quote}
Before the new one: say the Latin for bad faith, then the Latin for treacherous.
\end{quote}

\subsection*{You'll want to know: fallācia, fallāciae}
\textbf{fallācia, fallāciae} — \textquotedblleft{}a deceit, a trick\textquotedblright{}.

\begin{quote}
Fallāciīs eius nōn crēdō.
\end{quote}

\begin{quote}
\emph{I do not trust his tricks.}
\end{quote}

\emph{fallācia, fallāciae} f., built on \emph{fallō}, to deceive. A favourite word for a slave's trick in comedy. English \emph{fallacy}.

\begin{cousinweb}
\emph{fallācia} is built on \emph{fallāx} \textquotedblleft{}deceitful\textquotedblright{} (EtymDB), which comes from \emph{fallō} \textquotedblleft{}to deceive\textquotedblright{} (Lewis \& Short). English \emph{fallacy} comes from it.
\end{cousinweb}

\begin{grammarlens}[title={What you've built}]
One more word for telling, explaining and describing at B1.
\end{grammarlens}

\subsection*{Your turn}
\begin{itemize}
\raggedright
  \item \emph{Say it:} \emph{fallācia, fallāciae}
  \item \emph{Say it:} \emph{fallācia, fallāciae}, once more
  \item \emph{Say it:} a sentence of your own with \emph{fallācia, fallāciae}
  \item \emph{Recall:} say the Latin for bad faith, then the Latin for treacherous, then say \emph{fallācia, fallāciae} again
\end{itemize}

\begin{tcolorbox}[breakable,colback=teal!4,colframe=teal!35!black,title={Before you move on}]
What does fallācia, fallāciae mean? (\textquotedblleft{}A deceit, a trick\textquotedblright{}.) Say it once more.
\end{tcolorbox}
\section[dēcipiō, dēcipere]{dēcipiō, dēcipere — to deceive, to cheat}

\label{lesson:LA-C399-decipio}



\begin{quote}
Before the new one: say the Latin for treacherous, then the Latin for a deceit.
\end{quote}

\subsection*{You'll want to know: dēcipiō, dēcipere}
\textbf{dēcipiō, dēcipere} — \textquotedblleft{}to deceive, to cheat\textquotedblright{}.

\begin{quote}
Mercātor malus multōs dēcēpit.
\end{quote}

\begin{quote}
\emph{The dishonest merchant cheated many people.}
\end{quote}

\emph{dēcipiō, dēcipere, dēcēpī, dēceptum}, \emph{dē} + \emph{capiō}: to catch out. English \emph{deceive} and \emph{deception}.

\begin{cousinweb}
\emph{dēcipiō} is built from \emph{dē} and \emph{capiō} \textquotedblleft{}to take, to catch\textquotedblright{} (EtymDB; Lewis \& Short): to catch someone out. English \emph{deceive} comes from it through French.
\end{cousinweb}

\begin{grammarlens}[title={What you've built}]
One more word for telling, explaining and describing at B1.
\end{grammarlens}

\subsection*{Your turn}
\begin{itemize}
\raggedright
  \item \emph{Say it:} \emph{dēcipiō, dēcipere}
  \item \emph{Say it:} \emph{dēcipiō, dēcipere}, once more
  \item \emph{Say it:} a sentence of your own with \emph{dēcipiō, dēcipere}
  \item \emph{Recall:} say the Latin for treacherous, then the Latin for a deceit, then say \emph{dēcipiō, dēcipere} again
\end{itemize}

\begin{tcolorbox}[breakable,colback=teal!4,colframe=teal!35!black,title={Before you move on}]
What does dēcipiō, dēcipere mean? (\textquotedblleft{}To deceive, to cheat\textquotedblright{}.) Say it once more.
\end{tcolorbox}
\section[testor, testārī]{testor, testārī — to call to witness; to testify}

\label{lesson:LA-C399-testor}



\begin{quote}
Before the new one: say the Latin for a deceit, then the Latin for to deceive.
\end{quote}

\subsection*{You'll want to know: testor, testārī}
\textbf{testor, testārī} — \textquotedblleft{}to call to witness; to testify\textquotedblright{}.

\begin{quote}
Deōs testor mē vērum dīcere.
\end{quote}

\begin{quote}
\emph{I call the gods to witness that I speak the truth.}
\end{quote}

Deponent: \emph{testor, testārī, testātus sum}, built on \emph{testis}, a witness. \emph{Testāmentum}, a will, is its noun.

\begin{cousinweb}
\emph{testor} is built on \emph{testis} \textquotedblleft{}a witness\textquotedblright{} (EtymDB; Lewis \& Short). English \emph{testament} and \emph{testify} come from its family.
\end{cousinweb}

\begin{grammarlens}[title={What you've built}]
One more word for telling, explaining and describing at B1.
\end{grammarlens}

\subsection*{Your turn}
\begin{itemize}
\raggedright
  \item \emph{Say it:} \emph{testor, testārī}
  \item \emph{Say it:} \emph{testor, testārī}, once more
  \item \emph{Say it:} a sentence of your own with \emph{testor, testārī}
  \item \emph{Recall:} say the Latin for a deceit, then the Latin for to deceive, then say \emph{testor, testārī} again
\end{itemize}

\begin{tcolorbox}[breakable,colback=teal!4,colframe=teal!35!black,title={Before you move on}]
What does testor, testārī mean? (\textquotedblleft{}To call to witness; to testify\textquotedblright{}.) Say it once more.
\end{tcolorbox}
